%-------------------------------------------------------------------------------
%   CHAPTER 2 - ОНОЛЫН БОЛОН ХАРЬЦУУЛСАН СУДАЛГАА
%-------------------------------------------------------------------------------

\section{Онолын болон харьцуулсан судалгаа}
\label{sec:literature-review}

Энэ бүлэгт Монгол бичгийн түүх, flashcard сургалтын систем, орчин үеийн гар утасны аппликейшн хөгжүүлэлтийн технологиуд, мөн ижил төстэй хэл сурах аппликейшнуудын талаарх онолын судалгааг дэлгэрэнгүй авч үзнэ.

\subsection{Монгол бичгийн түүх ба онцлог}
\label{subsec:mongolian-script}

\subsubsection{Монгол бичгийн түүхэн хөгжил ба гүйцэтгэсэн үүрэг}

Монгол бичиг нь XIII зууны эхэн үед төрийн албан бичиг болж нэвтэрсэн. 1941 онд кирилл үсэгт шилжсэн хэдий ч 2020 онд төрийн албан хэрэгт хос бичгээр хөтлөх зорилт хэрэгжил эхлээд байна. Сургалтын технологийн шийдэл шаардлагатай.

\subsubsection{Монгол бичгийн бүтэц ба онцлог}

Монгол бичиг нь босоо бичгийн систем бөгөөд дээрээс доош, зүүнээс баруун тийш бичигддэг. Үндсэн онцлогууд:

\begin{itemize}
    \item \textbf{Үсгийн тоо:} 7 эгшиг, 26 гийгүүлэгч, нийт 33 үндсэн үсэг
    \item \textbf{Байрлалын хувилбарууд:} Нэг үсэг үгийн эхэнд, дунд, төгсгөлд өөр өөр хэлбэртэй бичигддэг
    \item \textbf{Холбох дүрэм:} Үсгүүд үгийн дотор холбогдож, үг бүрийг нэг бүтэн хэлбэрээр бичнэ
\end{itemize}

\paragraph{Сурахад тулгардаг хүндрэлүүд:}

\begin{itemize}
    \item Нэг үсгийн олон хэлбэр (3-4 өөр хэлбэр)
    \item Холбох дүрмийн нарийн ширийн зүйл
    \item Орчин үеийн сургалтын материалын хомсдол
\end{itemize}

\subsection{Мэдлэг эзэмших танин мэдэхүйн аргууд: Spaced Repetition болон Flashcard}
\label{subsec:flashcard-systems}

Шинэ үсэг, үг эзэмших үйл явцад Flashcard системүүд Ideavhtei сануулалт (Active Recall) болон Spaced Repetition зарчимд тулгуурлан үр дүнтэй ажилладаг~\cite{karpicke2008critical}. Hermannalbeau Ebbinghaus-ийн сэтгэл судлааны судалгаагаар давтах үйлдэлийн дараа мартуулах муруй буурч, мэдээлэл санамжид илүү удаан тогтдог болохыг батлажээ~\cite{ebbinghaus1985memory}. Anki болон Memrise зэрэг платформуудаас ашигладаг SuperMemo SM-2 алгоритм нь хувь хүний сур хүчнийн динамикт үндэслэн оптимал давтамжийг тооцоолдог~\cite{wozniak1990optimization}. Энэ онолоос Egüü системийг хөгжүүлэхэд нэмэлт мэдээлэл авахаа үе гэх мэт урьдчилгав заасан технологуудыг ашигласан.

\subsection{Гар утасны аппликейшн хөгжүүлэлтийн технологиуд}
\label{subsec:mobile-technologies}

\subsubsection{React Native - Cross-Platform фреймворк}

React Native нь Meta компаниас 2015 онд гаргасан JavaScript-ийн гар утасны фреймворк юм. iOS болон Android үйлдлийн системд нэгэн зэрэг ажиллах аппликейшн хөгжүүлэх боломжийг олгодог~\cite{nack2016learning}. Egüü төсөл хөгжүүлэхэд cross-platform дэмжлэг, Hot Reload, болон том коммьюнити ашиглалттай байсан давуу талууд.

\subsubsection{Firebase - Backend-as-a-Service}

Firebase нь Google-ийн үүлэн платформ бөгөөд хэрэглэгч бүртгэл (Authentication), NoSQL өгөгдлийн сан (Cloud Firestore), файл хадгалалт (Cloud Storage), болон Real-time synchronization үйлчилгээ олгодог~\cite{welekar2017firebase}. Serverbе инфраструктур удирдахгүй болгосон давуу талтай. Firestore өгөгдлийг collection-document бүтцээр хадгалж, Security Rules-ээр нарийвчилсан хандалтын хяналт хийдэг.

\subsection{TensorFlow.js ба OCR технологи}
\label{subsec:tensorflow-ocr}

TensorFlow.js нь Google-ийн машин сургалтын фреймворкийн JavaScript хувилбар бөгөөд төхөөрөмж дээр шууд ажилладаг (client-side inference) давуу талтай~\cite{smilkov2019tensorflow}. Хэрэглэгчийн өгөгдөл төхөөрөмжөөс гардаггүй, интернэтгүйгээр ч ажиллаж болно. OCR (Optical Character Recognition) системд хэрэглэгчийн гараар бичсэн Монгол бичгийг оруулаад зургийг preprocessing, segmentation, feature extraction, classification, post-processing дараалллээр боловсруулан текст болгож болно.

\subsection{Ижил төстэй аппликейшнуудын харьцуулсан судалгаа}
\label{subsec:comparative-analysis}

Хэл сурах аппликейшнуудаас Duolingo (gamification, 500 сая+ хэрэглэгч), Anki (SM-2 алгоритм, үр дүнтэй), Memrise (видео агуулга) гэх мэтүүд тэргүүлэгчүүд юм. Гэвч эдгээр аппуудаас Монгол бичиг дэмждэг, gamification + spaced repetition хослосон, орчин үеийн интерфэйстэй системүүд байхгүй байсан. Эдгээр цоорхойг нөхөхөөр Egüü төслийг хөгжүүлсэн.

\begin{table}[H]
\centering
\caption{Хэл сурах аппликейшнуудын харьцуулалт}
\label{tab:app-comparison}
\begin{tabular}{|l|c|c|c|c|}
\hline
\textbf{Онцлог} & \textbf{Duolingo} & \textbf{Anki} & \textbf{Memrise} & \textbf{Egüü} \\
\hline
Gamification & +++ & - & ++ & ++ \\
\hline
Spaced Repetition & + & +++ & ++ & +++ \\
\hline
UI/UX & +++ & + & ++ & +++ \\
\hline
Монгол бичиг & - & - & - & +++ \\
\hline
OCR таних & - & - & - & ++ \\
\hline
\end{tabular}
\end{table}
Egüü
\subsubsection{Egüü төслийн давуу талууд}

Дээрх судалгааны үндсэн дээр Egüü төсөл нь дараах давуу талуудтай:

\begin{enumerate}
    \item \textbf{Монгол бичигт тусгайлсан:} Одоогоор Монгол бичиг заах орчин үеийн апп байхгүй. Egüü нь энэ зах зээлийн цорын ганц шийдэл болно.
    
    \item \textbf{OCR технологи:} Хэрэглэгч гараар бичсэн үгээ таниулах боломжтой. Энэ нь бусад аппуудад байхгүй онцлог юм.
    
    \item \textbf{Орчин үеийн UI/UX:} React Native, NativeWind ашиглан орчин үеийн, сонирхолтой интерфэйс бүтээсэн.
    
    \item \textbf{Gamification + Spaced Repetition:} Duolingo-ийн gamification болон Anki-ийн spaced repetition-ийг хослуулсан.
    
    \item \textbf{Олон хэлний дэмжлэг:} i18n ашиглан Монгол, Англи хэлээр ашиглах боломжтой.
    
    \item \textbf{Үнэгүй, нээлттэй эх:} Бүх хүнд хүртээмжтэй, коммьюнитид хувь нэмэр оруулах боломжтой.
\end{enumerate}

\subsection{Дүгнэлт}

Энэхүү бүлэгт Монгол бичгийн түүх, flashcard сургалтын систем, орчин үеийн гар утасны аппликейшн хөгжүүлэлтийн технологиуд, мөн ижил төстэй хэл сурах аппликейшнуудын талаарх дэлгэрэнгүй судалгааг хийлээ.

Судалгааны үр дүнгээс харахад:
\begin{itemize}
    \item Монгол бичиг нь түүхэн өв болох уламжлалт бичиг боловч орчин үед сурах боломж хязгаарлагдмал байна
    \item Flashcard систем болон spaced repetition алгоритм нь хэл сурахад маш үр дүнтэй арга болох нь батлагдсан
    \item React Native, Firebase, TensorFlow.js зэрэг технологиуд нь орчин үеийн, үр дүнтэй гар утасны апп хөгжүүлэх боломжийг олгодог
    \item Одоогоор Монгол бичигт тусгайлсан орчин үеийн апп байхгүй тул Egüü төсөл нь чухал хэрэгцээг хангах болно
\end{itemize}

Дараагийн бүлэгт эдгээр онолын судалгаанд үндэслэн Egüü системийн шинжилгээ, зохиомжийг дэлгэрэнгүй танилцуулна.

\vfill